% !TeX root = ../../tesis.tex
% =============================================================================
%  5.2.1.3 FUERZA CAPILAR Y NODOS DE TRANSFERENCIA
% =============================================================================
\subsection{\bajoRevision{Fuerza Capilar y Nodos de Transferencia ($FC$)}}
\label{subsec:fuerza-capilar}
\notaRevision{Los resultados de esta subsección están sujetos a revisión con el
comité tutor. Los valores presentados corresponden a la corrida del modelo del
18 de abril de 2026.}

La \textbf{Fuerza Capilar} ($FC$) cuantifica el nivel de convergencia de flujos
de transporte en cada nodo o agrupación de nodos del grafo. El indicador se
evalúa desde dos perspectivas complementarias:

\textbf{A. Fuerza Capilar Geográfica ($FC_H$) — Macro-Hubs:} agrupa
topológicamente estaciones físicamente colindantes (radio de 100~m) en un
subgrafo $H$ que consolida los flujos de sus estaciones constituyentes,
representando el fenómeno urbano del transbordo físico (CETRAMs):

\begin{equation}
    FC_{H} = \sum_{j \in H} \bigl(k_{\text{in}}(j) + k_{\text{out}}(j)\bigr)
    - \text{Aristas Internas}
    \label{eq:fc-geografica}
\end{equation}

El Cuadro~\ref{tab:fuerza-capilar-hubs} detalla el comportamiento numérico de los
cinco principales Macro-Hubs de transferencia identificados en la red con base en
esta formulación, mientras que la Figura~\ref{fig:barplot-macro-hubs} permite
contrastar las magnitudes relativas de la concentración de flujos en los diez
nodos más significativos de la metrópoli.

\begin{table}[H]
    \centering
    \caption[Top 5 Macro-Hubs con mayor Fuerza Capilar Geográfica ($FC_H$)]{
    Los cinco Macro-Hubs de transferencia (nodos agrupados topológicamente a
    menos de 100~m) con mayor Fuerza Capilar Geográfica en la red metropolitana.
    Corrida: 2026-04-18.}
    \label{tab:fuerza-capilar-hubs}
    \fontsize{8.5pt}{10.5pt}\selectfont
    \begin{tabularx}{\textwidth}{|c|X|c|c|c|c|}
        \hline
        \textbf{ID Hub} & \textbf{Principales Estaciones Agrupadas} & \textbf{Sistemas} & \textbf{Nodos} & \textbf{Conexiones} \textit{(In / Out)} & \textbf{$FC_H$} \\
        \hline
        2174 & Zona Hospitales Sur (San Fernando, Chimalcoyotl, Coapa, Hosp. Manuel Gea Gtz., INER, Urgencias Respiratorias) & CC, RTP & 21 & 66 / 64 & \textbf{130} \\
        \hline
        136 & Eje Huipulco -- Estadio Azteca (Acoxpa, Huipulco, Estadio Azteca, Renato Leduc, San Alejandro) & CC, RTP, TL & 61 & 61 / 50 & \textbf{111} \\
        \hline
        1597 & Corredor Tlalpan Sur (Calle La Rosa, Tezoquipa, Galeana, Jojutla) & CC, RTP & 10 & 49 / 51 & \textbf{100} \\
        \hline
        259 & Nodo Barranca del Muerto (Barranca del Muerto, Alpes, Periférico Sur, José María Velasco, Miguel Ocaranza) & CC, RTP & 21 & 50 / 42 & \textbf{92} \\
        \hline
        129 & CETRAM Chapultepec (Metro Chapultepec, Acapulco, Agustín Melgar, Sonora, Veracruz, Tampico) & CC, METRO, RTP, TROLE & 29 & 37 / 40 & \textbf{77} \\
        \hline
    \end{tabularx}
    \vspace{1ex}
    \fuentefigura{Elaboración propia con \texttt{VFTModel} y datos de \texttt{Apimetro}.}
\end{table}

\begin{figure}[H]
    \centering
    \includegraphics[width=0.85\textwidth]{Figures/Cap5/FuerzaCapilar/Tabla/NB03_fig01_top10_macrohubs_barplot.png}
    \caption[Distribución de los principales Macro-Hubs por Fuerza Capilar]{
    Gráfico de barras comparativo de los diez principales Macro-Hubs de la red
    metropolitana, ordenados con base en su nivel acumulado de Fuerza Capilar
    Geográfica ($FC_H$).}
    \label{fig:barplot-macro-hubs}
    \fuentefigura{Elaboración propia con \texttt{VFTModel} a partir de los resultados analíticos del modelo.}
\end{figure}

Por último, la distribución espacial y la cobertura geográfica de estos nodos de
transferencia masiva dentro del entorno urbano de la \zmvm{} se presenta mapeada
de forma detallada en la Figura~\ref{fig:mapa-fuerza-capilar}.

\begin{figure}[H]
    \centering
    \includegraphics[width=0.85\textwidth]{Figures/Cap5/FuerzaCapilar/Mapa/NB03_mapa01_macrohubs.png}
    \caption[Mapa de Macro-Hubs y Fuerza Capilar Geográfica]{
    Representación geoespacial de la \termino{Fuerza Capilar Geográfica}
    ($FC_H$). El mapa destaca la concentración de flujos de transporte en los
    principales nodos y zonas de transferencia física (Macro-Hubs) de la \zmvm.}
    \label{fig:mapa-fuerza-capilar}
    \fuentefigura{Elaboración propia con \texttt{VFTModel} utilizando visualización geoespacial.}
\end{figure}

% =============================================================================
%  5.2.1.4 FUERZA CAPILAR SISTÉMICA (MATEMÁTICA)
% =============================================================================
\subsection{\bajoRevision{Fuerza Capilar Sistémica (Matemática)}}
\label{subsec:fuerza-capilar-matematica}

\textbf{B. Fuerza Capilar Sistémica o Matemática ($FC$):} A diferencia de la
aproximación geográfica orientada a macro-hubs, esta perspectiva evalúa de manera
estricta la convergencia de flujos directamente sobre cada nodo topológico
individual. El análisis permite identificar los puntos exactos de la
infraestructura matemática de la red donde se concentra la mayor densidad de
aristas de entrada y salida, aislando singularidades críticas en la traza sin
aplicar tolerancias de proximidad espacial ni fusión de estaciones colindantes.

El Cuadro~\ref{tab:fuerza-capilar-matematica} detalla el comportamiento de los
cinco nodos individuales que registraron la mayor Fuerza Capilar pura dentro del
grafo de transporte metropolitano.

\begin{table}[H]
    \centering
    \caption[Top 5 Nodos con mayor Fuerza Capilar Sistémica (Matemática)]{
    Los cinco nodos individuales con mayor Fuerza Capilar Sistémica (Matemática)
    en la red metropolitana, evaluados de manera estricta sin agrupación
    geográfica. Corrida: 2026-04-18.}
    \label{tab:fuerza-capilar-matematica}
    \fontsize{9.5pt}{11.5pt}\selectfont
    \begin{tabular}{|c|l|c|c|c|}
        \hline
        \textbf{ID Nodo} & \textbf{Estación (Nodo Topológico)} & \textbf{Sistemas} & \textbf{Conexiones} \textit{(In / Out)} & \textbf{$FC$} \\
        \hline
        7270  & Luis Murillo                          & RTP & 21 / 20 & \textbf{41} \\
        \hline
        2817  & Calz. de Tlalpan -- Glorieta Huipulco & RTP & 19 / 15 & \textbf{34} \\
        \hline
        10211 & Sobre Calz. de Tlalpan y Luis Murillo & CC  & 16 / 17 & \textbf{33} \\
        \hline
        10717 & Tlalpan y Luis Murillo                & CC  & 18 / 15 & \textbf{33} \\
        \hline
        5540  & Estadio Azteca 2                      & RTP & 14 / 18 & \textbf{32} \\
        \hline
    \end{tabular}
    \vspace{1ex}
    \fuentefigura{Elaboración propia con \texttt{VFTModel} y datos de \texttt{Apimetro}.}
\end{table}

Las magnitudes y diferencias relativas entre los diez elementos individuales con
mayor concentración de conexiones se visualizan en el gráfico de barras de la
Figura~\ref{fig:barplot-nodos-matematicos}.

\begin{figure}[H]
    \centering
    \includegraphics[width=0.85\textwidth]{Figures/Cap5/FuerzaCapilar/Tabla/NB03_fig02_top10_nodos_matematicos_barplot.png}
    \caption[Distribución de los principales nodos individuales por Fuerza Capilar]{
    Gráfico de barras comparativo de los diez nodos individuales con mayor
    Fuerza Capilar Sistémica (Matemática) en la red metropolitana.}
    \label{fig:barplot-nodos-matematicos}
    \fuentefigura{Elaboración propia con \texttt{VFTModel} a partir de los resultados analíticos del modelo.}
\end{figure}

La distribución geográfica e implicación espacial de estos puntos específicos de
alta carga conectiva dentro de la red urbana consolidada se presenta en la
Figura~\ref{fig:mapa-nodos-matematicos}.

\begin{figure}[H]
    \centering
    \includegraphics[width=0.85\textwidth]{Figures/Cap5/FuerzaCapilar/Mapa/NB03_mapa02_nodos_matematicos.png}
    \caption[Mapa de Nodos Individuales y Fuerza Capilar Sistémica]{
    Representación geoespacial de la Fuerza Capilar Sistémica (Matemática) a
    nivel de nodo individual en la \zmvm{}.}
    \label{fig:mapa-nodos-matematicos}
    \fuentefigura{Elaboración propia con \texttt{VFTModel} utilizando visualización geoespacial.}
\end{figure}

\subsubsection{Análisis del Eje Periférico — Infraestructura Preexistente}
\label{subsubsec:periferico-fc}

Como diagnóstico previo a la propuesta del corredor anillar, se aislaron los nodos
y Macro-Hubs cuyos nombres contienen la denominación ``Periférico'', con el objetivo
de identificar los puntos de mayor tensión vehicular e intermodal sobre esta
vialidad fundamental. Este análisis se divide en dos enfoques: la concentración
geográfica (transbordos físicos en Macro-Hubs) y la concentración matemática
estricta (nodos topológicos individuales).

\paragraph{A. Fuerza Capilar Geográfica (Macro-Hubs) en el Periférico}

Al aplicar un radio de tolerancia de 100~m para agrupar las estaciones colindantes
sobre el eje, emergen los principales complejos de transbordo y alimentación. El
Cuadro~\ref{tab:periferico-hubs-geo} detalla los cinco Macro-Hubs de mayor Fuerza
Capilar Geográfica ($FC_H$), destacando la importancia de Barranca del Muerto y la
intersección con la zona de hospitales/Tlalpan como los anclajes más fuertes del
corredor.

\begin{table}[H]
    \centering
    \caption[Top 5 Macro-Hubs con mayor Fuerza Capilar sobre el Eje Periférico]{
    Cinco Macro-Hubs de transferencia con mayor Fuerza Capilar Geográfica
    ($FC_H$) ubicados sobre el Anillo Periférico. Corrida: 2026-04-18.}
    \label{tab:periferico-hubs-geo}
    \fontsize{8.5pt}{10.5pt}\selectfont
    \begin{tabularx}{\textwidth}{|c|X|c|c|c|c|}
        \hline
        \textbf{ID Hub} & \textbf{Principales Estaciones Agrupadas} & \textbf{Sistemas} & \textbf{Nodos} & \textbf{Conexiones} \textit{(In / Out)} & \textbf{$FC_H$} \\
        \hline
        259 & Nodo Barranca del Muerto (Barranca del Muerto, Alpes, Periférico Sur, Miguel Ocaranza) & CC, RTP & 21 & 50 / 42 & \textbf{92} \\
        \hline
        374 & Intersección Tlalpan--Periférico (Renato Leduc, Huipulco, Periférico, Av. Pedregal) & CC, RTP & 18 & 34 / 33 & \textbf{67} \\
        \hline
        453 & Nodo San Antonio (San Antonio, Periférico Sur, Puente Periférico Norte) & CC, RTP & 7 & 23 / 23 & \textbf{46} \\
        \hline
        402 & Eje Oriente Tepalcates (Constitución de Apatzingán, Canal de San Juan, Tepalcates) & CC, MB, RTP & 8 & 20 / 26 & \textbf{46} \\
        \hline
        291 & Intersección Xochimilco--Tepepan (16 de Septiembre, Kioto, Tec de Monterrey, Periférico Participación) & CC, RTP, TL & 14 & 23 / 22 & \textbf{45} \\
        \hline
    \end{tabularx}
    \vspace{1ex}
    \fuentefigura{Elaboración propia con \texttt{VFTModel} y datos de \texttt{Apimetro}.}
\end{table}

La Figura~\ref{fig:barplot-hubs-periferico} compara las magnitudes relativas de
estos agrupamientos geográficos, mientras que la Figura~\ref{fig:mapa-hubs-periferico}
espacializa su huella sobre la vialidad, revelando los polos naturales de atracción
de demanda que la futura infraestructura anillar deberá integrar.

\begin{figure}[H]
    \centering
    \includegraphics[width=0.85\textwidth]{Figures/Cap5/FuerzaCapilar/Tabla/NB03_fig04_top10_macrohubs_periferico_barplot.png}
    \caption[Fuerza Capilar de Macro-Hubs en Anillo Periférico]{Gráfico de barras
    de los diez principales Macro-Hubs geográficos ubicados en el Anillo Periférico,
    ordenados por su nivel de Fuerza Capilar.}
    \label{fig:barplot-hubs-periferico}
    \fuentefigura{Elaboración propia con \texttt{VFTModel}.}
\end{figure}

\begin{figure}[H]
    \centering
    \includegraphics[width=0.85\textwidth]{Figures/Cap5/FuerzaCapilar/Mapa/NB03_mapa03_macrohubs_periferico.png}
    \caption[Distribución espacial de Macro-Hubs en el Periférico]{Mapa de la
    Fuerza Capilar Geográfica ($FC_H$) en los Macro-Hubs del Anillo Periférico.}
    \label{fig:mapa-hubs-periferico}
    \fuentefigura{Elaboración propia con \texttt{VFTModel} utilizando visualización geoespacial.}
\end{figure}

\paragraph{B. Fuerza Capilar Sistémica (Nodos Matemáticos) en el Periférico}

Al desagregar los Hubs y analizar la Fuerza Capilar Sistémica ($FC$) estricta a
nivel de nodo individual, es posible identificar los paraderos específicos que
soportan la mayor carga de conexiones en la infraestructura actual. El
Cuadro~\ref{tab:periferico-nodos-math} presenta los cinco paraderos matemáticos
de mayor convergencia de flujo.

\begin{table}[H]
    \centering
    \caption[Nodos con mayor Fuerza Capilar Sistémica sobre el Eje Periférico]{
    Cinco nodos topológicos individuales con mayor Fuerza Capilar Matemática
    ($FC$) ubicados sobre el Anillo Periférico. Corrida: 2026-04-18.}
    \label{tab:periferico-nodos-math}
    \fontsize{9.5pt}{11.5pt}\selectfont
    \begin{tabular}{|c|l|c|r|r|c|}
        \hline
        \textbf{ID Nodo} & \textbf{Estación (Nodo Topológico)} & \textbf{Sistema} & \textbf{Entr.} & \textbf{Sal.} & \textbf{$FC$} \\
        \hline
        707  & Anillo Periférico y Blvd. Adolfo Ruiz Cortines & RTP & 14 & 16 & \textbf{30} \\
        \hline
        706  & Anillo Periférico y Blvd. Adolfo Ruiz Cortines & RTP & 13 & 13 & \textbf{26} \\
        \hline
        8538 & Periférico -- Barranca del Muerto              &  CC & 12 &  9 & \textbf{21} \\
        \hline
        602  & Anillo Periférico Canal de Garay -- Bilbao     & RTP & 10 & 10 & \textbf{20} \\
        \hline
        8510 & Periférico                                     & RTP &  9 & 11 & \textbf{20} \\
        \hline
    \end{tabular}
    \vspace{1ex}
    \fuentefigura{Elaboración propia con \texttt{VFTModel} y datos de \texttt{Apimetro}.}
\end{table}

La jerarquía analítica de estos puntos singulares se constata en la
Figura~\ref{fig:barplot-nodos-periferico}, seguida de la
Figura~\ref{fig:mapa-nodos-periferico}, que detalla la distribución geográfica a
escala de paradero individual, útil para el micro-posicionamiento de futuras
estaciones anillares.

\begin{figure}[H]
    \centering
    \includegraphics[width=0.85\textwidth]{Figures/Cap5/FuerzaCapilar/Tabla/NB03_fig05_top10_nodos_matematicos_periferico_barplot.png}
    \caption[Fuerza Capilar de nodos individuales en Anillo Periférico]{
    Gráfico de barras comparativo de los diez principales nodos matemáticos
    individuales ubicados en el Anillo Periférico.}
    \label{fig:barplot-nodos-periferico}
    \fuentefigura{Elaboración propia con \texttt{VFTModel}.}
\end{figure}

\begin{figure}[H]
    \centering
    \includegraphics[width=0.85\textwidth]{Figures/Cap5/FuerzaCapilar/Mapa/NB03_mapa04_nodos_matematicos_periferico.png}
    \caption[Distribución espacial de nodos matemáticos en el Periférico]{Mapa
    de la Fuerza Capilar Sistémica (Matemática) a nivel de nodo individual sobre
    el corredor del Anillo Periférico.}
    \label{fig:mapa-nodos-periferico}
    \fuentefigura{Elaboración propia con \texttt{VFTModel} utilizando visualización geoespacial.}
\end{figure}

\vspace{2ex}
\begin{quote}
    \small\textit{Nota aclaratoria:} los datos presentados en esta sección
    reflejan exclusivamente la infraestructura de transporte \textbf{preexistente}
    sobre el Anillo Periférico. No se ha trazado ni modelado aún la geometría
    del corredor anillar propuesto; la sección permanece en plano especulativo
    para identificar los puntos naturales de anclaje de la futura línea a nivel
    capilar.
\end{quote}
